\documentclass[11pt]{amsart}
\usepackage{amssymb,amsmath}
\newtheorem{theorem}{Theorem}
\newtheorem{proposition}[theorem]{Proposition}
\newtheorem{corollary}[theorem]{Corollary}
\newtheorem{lemma}[theorem]{Lemma}
\theoremstyle{remark}
\newtheorem{remark}[theorem]{Remark}

\title[The pentagonal bipyramid is a strict local minimum]{The pentagonal bipyramid is a strict local minimum of the logarithmic energy of seven points on the sphere}
\author{Draft: not reviewed}
\date{26 September 2026 (draft)}

\begin{document}
\begin{abstract}
The Hessian of the logarithmic energy of seven points on the unit sphere is degenerate at the
pentagonal bipyramid: besides the three rotations it has a two-dimensional kernel, so second
derivatives decide neither local minimality nor the stability of the corresponding equilibrium of
seven identical point vortices (Constantineau, Garc\'{\i}a-Azpeitia, Garc\'{\i}a-Naranjo and Lessard,
Remark~6.3). We show by a computer-assisted proof that the bipyramid is nevertheless a strict local
minimum modulo rotations, with an explicit neighbourhood, and hence a Lyapunov stable equilibrium
modulo rotations. After the nondegenerate directions are eliminated, the energy on the degenerate plane
grows exactly like one tenth of the fourth power of the displacement. We also show that a
three-point certificate of the type used by Kryvonos, Liehr and Taylor for eight points cannot prove
global optimality for seven points unless its minorant has fourth-order contact at the inner product $0$.
\end{abstract}
\maketitle

\section{Statement}
For $X=(x_0,\dots,x_6)\in(S^2)^7$ let $E(X)=-\sum_{i<j}\log|x_i-x_j|$. Let $B$ be the labelled
bipyramid: $b_0=(0,0,1)$, $b_1=(0,0,-1)$, $b_{2+k}=(\cos\frac{2\pi k}{5},\sin\frac{2\pi k}{5},0)$. Then
$E(B)=-6\log 2-\frac52\log 5$.

\begin{theorem}\label{thm:local}
If $g\in SO(3)$ and $|x_i-g b_i|\le 7.2\cdot 10^{-5}$ for all $i$, and $X$ is not a rotation of $B$, then $E(X)>E(B)$.
\end{theorem}

\begin{theorem}\label{thm:stab}
The bipyramid is an equilibrium of seven identical point vortices on the sphere, and it is Lyapunov
stable modulo rotations: for every $\varepsilon>0$ there is $\delta>0$ with
$\operatorname{dist}(X(0),SO(3)B)<\delta\Rightarrow\operatorname{dist}(X(t),SO(3)B)<\varepsilon$ for all $t$.
\end{theorem}

Theorem~\ref{thm:stab} follows from Theorem~\ref{thm:local} because the Hamiltonian
$H=-\sum_{i<j}\log|x_i-x_j|^2=2E$ is conserved and the orbit $SO(3)B$ is compact: a strict minimum on a
punctured neighbourhood of a compact invariant set makes its small sublevel sets invariant neighbourhoods.

\section{Chart}
Rotate $x_0$ to the north pole and project the other points stereographically from it,
$z=(X+iY)/(1-Z)$. Then
\[
E=F(z):=-\sum_{1\le i<j\le 6}\log|z_i-z_j|+3\sum_{i=1}^6\log(1+|z_i|^2)-21\log 2,
\]
with $B$ at $z_S=0$, $z_k=\omega^k$ ($\omega=e^{2\pi i/5}$). Coordinates $w\in\mathbb R^{11}$:
$z_S=x+iy$, $z_k=\omega^k(1+u_k+iv_k)$, $v_0=0$. All data at $w=0$ lie in $K=\mathbb Q(\sin\frac{2\pi}5)$.

\section{The degenerate expansion (exact)}
In exact arithmetic in $K$: $\nabla F(0)=0$; the Hessian $H$ has rank $9$; with $N$ a basis of $\ker H$
and $C$ a basis of $\operatorname{range}H$ (exact, $C^TN=0$), $A=C^THC\succ 0.2659\,I$; the cubic vanishes
on $\ker H$; with $b(\xi)=C^T\nabla F_3(N\xi)$ and $\eta^*(\xi)=-A^{-1}b(\xi)$,
\[
q(\xi)=F_4(N\xi)-\tfrac12 b(\xi)^TA^{-1}b(\xi)=\tfrac1{10}|N\xi|^4 .
\]
The kernel moves only the pentagon, in latitude, in the wavenumber-two pattern
$u_k=a\cos\frac{4\pi k}5+b\sin\frac{4\pi k}5$.

\section{Proof of Theorem~\ref{thm:local}}
Put $G(\xi,\zeta)=F(N\xi+C(\eta^*(\xi)+\zeta))-F(0)$. Rigorous ball arithmetic gives, for
$|\xi|\le 10^{-3}$: $G(\xi,0)\ge 0.0994|\xi|^4$ (power series along rays, enclosed over arcs by a
mean-value form, plus a Cauchy tail on the disc $|r|\le 0.12$);
$|\nabla_\zeta G(\xi,0)|\le 2.979|\xi|^3$ (same method); and, on a box containing the domain with
$|\zeta|\le 2\cdot10^{-3}$, $\lambda_{\min}(C^T\nabla^2F\,C)\ge\mu=0.0664$ (third-derivative enclosure and
Weyl's inequality). Taylor's theorem in $\zeta$ and $2.979\,X^3Y\le\frac{\mu}{4}Y^2+\frac{2.979^2}{\mu}X^6$ give
\[
G(\xi,\zeta)\ge 0.0166\,|\zeta|^2+0.0992\,|\xi|^4 ,
\]
which is positive unless $\xi=\zeta=0$. The chart ball $|w|\le 6.6\cdot10^{-4}$ lies in the domain, and a
configuration within $7.2\cdot10^{-5}$ of a rotated $B$ has a rotation in that ball (explicit rotation and
stereographic estimates). The programs and all constants are in the accompanying folder.

\section{An obstruction for the global problem}
\begin{proposition}
Let $H\le\phi$ on $[-1,1)$, $\phi(t)=-\frac12\log(2-2t)$, with $H(0)=\phi(0)$ and $(\phi-H)''(0)>0$.
Then $\inf_Y\sum_{i<j}H(y_i\cdot y_j)<E(B)$.
\end{proposition}
\begin{proof}
Along $w(s)=sN\xi_0$ the energy satisfies $E(Y(s))-E(B)=O(s^4)$, while each pole--pentagon inner product
is $\pm u_k(s)+O(s^2)$ with $u_k(s)=s(N\xi_0)_{u_k}\not\equiv0$. Hence
$\sum_{i<j}H(y_i\cdot y_j)\le E(Y(s))-c\,s^2+O(s^3)<E(B)$ for small $s\neq0$.
\end{proof}
Kryvonos, Liehr and Taylor's minorants have exactly double contact at interior nodes; for seven points a
sharp certificate would need contact of order four at $t=0$.

\begin{remark}
Global optimality of the bipyramid remains open. By Armentano et al., it would follow from showing that
some minimizer contains an antipodal pair.
\end{remark}

\begin{thebibliography}{9}
\bibitem{CGGL} K.~Constantineau, C.~Garc\'{\i}a-Azpeitia, L.~C.~Garc\'{\i}a-Naranjo, J.-P.~Lessard,
\emph{Determination of stable branches of relative equilibria of the $N$-vortex problem on the sphere},
arXiv:2309.04320.
\bibitem{KLT} L.~Kryvonos, L.~Liehr, M.~A.~Taylor, \emph{Energy minimization for eight points on the sphere},
arXiv:2609.22077.
\bibitem{ABCFVV} D.~Armentano, L.~Bentancur, F.~Carrasco, M.~Fiori, M.~Vald\'es, M.~Velasco,
\emph{Characterization of logarithmic Fekete critical configurations of at most six points in all dimensions},
arXiv:2502.10152.
\bibitem{CW} H.~Cohn, J.~Woo, \emph{Three-point bounds for energy minimization}, J. Amer. Math. Soc. 25 (2012).
\bibitem{B} C.~Beltr\'an, \emph{On Smale's 7th problem} (Spanish), Gac. R. Soc. Mat. Esp. 23 (2020).
\bibitem{Arb} F.~Johansson, \emph{Arb: efficient arbitrary-precision midpoint-radius interval arithmetic},
IEEE Trans. Comput. 66 (2017).
\end{thebibliography}
\end{document}
